\begin{tcolorbox}[colback=cadmiumgreen!5!white,colframe=forestgreen!75!white,breakable,title=\textit{Optimized setting prompt template for the mortality prediction task on the TJH.}]
\begin{VerbatimWrap}
You are an experienced doctor specializing in COVID-19 treatment, skilled in interpreting longitudinal patient data and predicting clinical outcomes.

I will provide you with longitudinal medical information for a patient. The data covers 5 visits/time points that occurred at 2020-01-23, 2020-01-30, 2020-02-04, 2020-02-05, 2020-02-06.
Each clinical feature is presented as a list of values, corresponding to these time points. Missing values are represented as `NaN` for numerical values and "unknown" for categorical values. Note that units and reference ranges are provided alongside relevant features.

Patient Background:
- Sex: female
- Age: 70.0 years

Your Task:
Your primary task is to assess the provided medical data and analyze the health records from ICU visits to determine the likelihood of the patient not surviving their hospital stay.

Instructions & Output Format:
Please first perform a step-by-step analysis of the patient data, considering trends, abnormal values relative to reference ranges, and their clinical significance for survival. Then, provide a final assessment of the likelihood of not surviving the hospital stay.

Your final output must be a JSON object containing two keys:
1.  `"think"`: A string containing your detailed step-by-step clinical reasoning (under 500 words).
2.  `"answer"`: A floating-point number between 0 and 1 representing the predicted probability of mortality (higher value means higher likelihood of death).

Example Format: ```json { "think": "The patient presents with worsening X, stable Y, and improved Z. Factor A is a major risk indicator... Overall assessment suggests a high risk.", "answer": 0.85 }```

Handling Uncertainty:
In situations where the provided data is clearly insufficient or too ambiguous to make a reasonable prediction, respond with the exact phrase: `I do not know`

Now, please analyze and predict for the following patient:

Clinical Features Over Time:
- Hypersensitive cardiac troponinI (Unit: ng/L. Reference range: less than 14.): [NaN, NaN, NaN, NaN, NaN]
- hemoglobin (Unit: g/L. Reference range: 140 - 180 for men, 120 - 160 for women.): [109.0, 112.0, NaN, 126.0, NaN]
- Serum chloride (Unit: mmol/L. Reference range: 96 - 106.): [99.1, 102.9, NaN, 102.2, NaN]
- Prothrombin time (Unit: seconds. Reference range: 13.1 - 14.125.): [13.6, NaN, NaN, NaN, NaN]
- procalcitonin (Unit: ng/mL. Reference range: less than 0.05.): [0.06, 0.06, NaN, NaN, NaN]
- eosinophils(%) (Unit: %. Reference range: 1 - 6.): [0.0, 0.2, NaN, 0.1, NaN]
- Interleukin 2 receptor (Unit: pg/mL. Reference range: less than 625.): [591.0, NaN, NaN, NaN, NaN]
- Alkaline phosphatase (Unit: IU/L. Reference range: 44 - 147.): [47.0, 61.0, NaN, 69.0, NaN]
- albumin (Unit: g/dL. Reference range: 3.5 - 5.5.): [34.9, 34.0, NaN, 38.4, NaN]
- basophil(%) (Unit: %. Reference range: 0.5 - 1.): [0.0, 0.2, NaN, 0.1, NaN]
- Interleukin 10 (Unit: pg/mL. Reference range: less than 9.8.): [8.1, NaN, NaN, NaN, NaN]
- Total bilirubin (Unit: µmol/L. Reference range: 5.1 - 17.): [7.5, 7.7, NaN, 12.6, NaN]
- Platelet count (Unit: x 10^9/L. Reference range: 150 - 450.): [169.0, 275.0, NaN, 238.0, NaN]
- monocytes(%) (Unit: %. Reference range: 2 - 10.): [5.9, 6.4, NaN, 6.1, NaN]
- antithrombin (Unit: %. Reference range: 80 - 120.): [84.0, NaN, NaN, NaN, NaN]
- Interleukin 8 (Unit: pg/mL. Reference range: less than 62.): [21.9, NaN, NaN, NaN, NaN]
- indirect bilirubin (Unit: μmol/L. Reference range: 3.4 - 12.0.): [3.8, 3.7, NaN, 7.5, NaN]
- Red blood cell distribution width (Unit: %. Reference range: 11.5 - 14.5 for men, 12.2 - 16.1 for women.): [12.8, 12.6, NaN, NaN, NaN]
- neutrophils(%) (Unit: %. Reference range: 45 - 70.): [75.0, 60.1, NaN, 66.4, NaN]
- total protein (Unit: g/L. Reference range: 60 - 83.): [68.3, 62.2, NaN, 68.2, NaN]
- Quantification of Treponema pallidum antibodies (Unit: /. Reference range: less than 1.0.): [0.07, NaN, NaN, NaN, NaN]
- Prothrombin activity (Unit: %. Reference range: 70 - 130.): [94.0, NaN, NaN, NaN, NaN]
- HBsAg (Unit: IU/mL. Reference range: 0.0 - 0.01.): [0.01, NaN, NaN, NaN, NaN]
- mean corpuscular volume (Unit: fL. Reference range: 80 - 100.): [94.6, 93.6, NaN, 94.4, NaN]
- hematocrit (Unit: %. Reference range: 40 - 54 for men, 36 - 48 for women.): [33.2, 33.7, NaN, 35.7, NaN]
- White blood cell count (Unit: x 10^9/L. Reference range: 4.5 - 11.0.): [3.72, 5.31, NaN, 7.68, NaN]
- Tumor necrosis factorα (Unit: pg/mL. Reference range: less than 8.1.): [10.4, NaN, NaN, NaN, NaN]
- mean corpuscular hemoglobin concentration (Unit: g/L. Reference range: 320 - 360.): [328.0, 332.0, NaN, 353.0, NaN]
- fibrinogen (Unit: g/L. Reference range: 2 - 4.): [5.33, NaN, NaN, NaN, NaN]
- Interleukin 1β (Unit: pg/mL. Reference range: less than 6.5.): [5.0, NaN, NaN, NaN, NaN]
- Urea (Unit: mmol/L. Reference range: 1.8 - 7.1.): [2.9, 3.7, NaN, 4.22, NaN]
- lymphocyte count (Unit: x 10^9/L. Reference range: 1.0 - 4.8.): [0.71, 1.76, NaN, 2.1, NaN]
- PH value (Unit: /. Reference range: 7.35 - 7.45.): [NaN, NaN, NaN, NaN, NaN]
- Red blood cell count (Unit: x 10^12/L. Reference range: 4.5 - 5.5 for men, 4.0 - 5.0 for women.): [3.51, 3.6, NaN, 3.78, NaN]
- Eosinophil count (Unit: x 10^9/L. Reference range: 0.02 - 0.5.): [0.0, 0.01, NaN, 0.01, NaN]
- Corrected calcium (Unit: mmol/L. Reference range: 2.12 - 2.57.): [2.17, 2.25, NaN, 2.32, NaN]
- Serum potassium (Unit: mmol/L. Reference range: 3.5 - 5.0.): [3.34, 3.34, NaN, 3.9, NaN]
- glucose (Unit: mmol/L. Reference range: 3.9 - 5.6.): [9.02, 9.42, NaN, NaN, NaN]
- neutrophils count (Unit: x 10^9/L. Reference range: 2.0 - 8.0.): [2.79, 3.19, NaN, 5.09, NaN]
- Direct bilirubin (Unit: µmol/L. Reference range: 1.7 - 5.1.): [3.7, 4.0, NaN, 5.1, NaN]
- Mean platelet volume (Unit: fL. Reference range: 7.4 - 11.4.): [11.6, 10.7, NaN, 9.7, NaN]
- ferritin (Unit: ng/mL. Reference range: 24 - 336 for men, 11 - 307 for women.): [567.2, NaN, NaN, NaN, NaN]
- RBC distribution width SD (Unit: fL. Reference range: 40.0 - 55.0.): [44.4, 42.7, NaN, NaN, NaN]
- Thrombin time (Unit: seconds. Reference range: 12 - 19.): [16.7, NaN, NaN, NaN, NaN]
- (%)lymphocyte (Unit: %. Reference range: 20 - 40.): [19.1, 33.1, NaN, 27.3, NaN]
- HCV antibody quantification (Unit: IU/mL. Reference range: 0.04 - 0.08.): [0.06, NaN, NaN, NaN, NaN]
- DD dimer (Unit: mg/L. Reference range: 0 - 0.5.): [0.98, NaN, NaN, NaN, NaN]
- Total cholesterol (Unit: mmol/L. Reference range: less than 5.17.): [3.28, 3.49, NaN, 4.47, NaN]
- aspartate aminotransferase (Unit: U/L. Reference range: 8 - 33.): [30.0, 30.0, NaN, 20.0, NaN]
- Uric acid (Unit: µmol/L. Reference range: 240 - 510 for men, 160 - 430 for women.): [151.0, 135.0, NaN, 221.1, NaN]
- HCO3 (Unit: mmol/L. Reference range: 22 - 29.): [22.3, 24.6, NaN, 25.3, NaN]
- calcium (Unit: mmol/L. Reference range: 2.13 - 2.55.): [2.07, 2.13, NaN, 2.29, NaN]
- Aminoterminal brain natriuretic peptide precursor(NTproBNP) (Unit: pg/mL. Reference range: 0 - 125.): [NaN, NaN, NaN, NaN, NaN]
- Lactate dehydrogenase (Unit: U/L. Reference range: 140 - 280.): [328.0, 269.0, NaN, 226.0, NaN]
- platelet large cell ratio (Unit: %. Reference range: 15 - 35.): [37.8, 30.0, NaN, 21.4, NaN]
- Interleukin 6 (Unit: pg/mL. Reference range: 0 - 7.): [47.82, NaN, NaN, NaN, NaN]
- Fibrin degradation products (Unit: μg/mL. Reference range: 0 - 10.): [4.0, NaN, NaN, NaN, NaN]
- monocytes count (Unit: x 10^9/L. Reference range: 0.32 - 0.58.): [0.22, 0.34, NaN, 0.47, NaN]
- PLT distribution width (Unit: fL. Reference range: 9.2 - 16.7.): [13.9, 12.5, NaN, 10.1, NaN]
- globulin (Unit: g/L. Reference range: 23 - 35.): [33.4, 28.2, NaN, 29.8, NaN]
- γglutamyl transpeptidase (Unit: U/L. Reference range: 7 - 47 for men, 5 - 25 for women.): [21.0, 54.0, NaN, 53.0, NaN]
- International standard ratio (Unit: ratio. Reference range: 0.8 - 1.2.): [1.04, NaN, NaN, NaN, NaN]
- basophil count(#) (Unit: x 10^9/L. Reference range: 0.01 - 0.02.): [0.0, 0.01, NaN, 0.01, NaN]
- mean corpuscular hemoglobin (Unit: pg. Reference range: 27 - 31.): [31.1, 31.1, NaN, 33.3, NaN]
- Activation of partial thromboplastin time (Unit: seconds. Reference range: 22 - 35.): [34.8, NaN, NaN, NaN, NaN]
- High sensitivity Creactive protein (Unit: mg/L. Reference range: 3 - 10.): [42.3, 3.6, NaN, NaN, NaN]
- HIV antibody quantification (Unit: IU/mL. Reference range: 0.08 - 0.11.): [0.1, NaN, NaN, NaN, NaN]
- serum sodium (Unit: mmol/L. Reference range: 135 - 145.): [135.7, 140.4, NaN, 143.2, NaN]
- thrombocytocrit (Unit: %. Reference range: 0.22 - 0.24.): [0.2, 0.29, NaN, 0.23, NaN]
- ESR (Unit: mm/hr. Reference range: less than 15  for men, less than 20 for women.): [66.0, 29.0, NaN, NaN, NaN]
- glutamicpyruvic transaminase (Unit: U/L. Reference range: 0 - 35.): [19.0, 84.0, NaN, 67.0, NaN]
- eGFR (Unit: mL/min/1.73m². Reference range: more than 90.): [77.2, 90.0, NaN, 84.6, NaN]
- creatinine (Unit: µmol/L. Reference range: 61.9 - 114.9 for men, 53 - 97.2 for women.): [69.0, 58.0, NaN, 64.0, NaN]
\end{VerbatimWrap}
\end{tcolorbox}